% ============================================================
% Discussion — Paper A (JNCA): interpretation, implications, limitations
% ============================================================
\section{Discussion}
\label{sec:discussion}

The cascade invariant addresses the failure mode that sending-side defenses cannot cover, and the rigor
analysis marks where the signal stops being operationally useful. We focus here on the mechanism's
implications and its limits.

\subsection{Interpretation: why the receiving side degrades less}
A volumetric detector in the spirit of COMPA \citep{stringhini2015thataint} assumes that a compromise shows up
as a surge; a path-based detector in the spirit of Hopper \citep{ho2021hopper} assumes an unusual route to a
rare target. Both read what the attacker controls. When the attacker deliberately spreads the sends over
time, the volumetric signal vanishes and its AUC drops below the random baseline, and the path-based signal
degrades as the cascade reaches further. Cascade dynamics instead respond to the recipient side, the
wave-like signature among the neighbors, which the sender cannot hide without giving up fan-out
altogether---and a single recipient is exactly the surgical regime, which we mark as the blind spot. Pacing
the sends weakens this signal (AUC $0.60$ at $K{=}1$) but much less than it weakens volume. This
agrees with Ho et al.\ \citep{ho2019detecting} that lateral phishing needs a signal beyond the content of a
single message; we differ in locating that signal in the dynamics of the communication graph rather than in
headers or relays. The time-shuffle control supports the temporal reading: shuffling removes the temporal
order (and, as a side effect, within-bucket co-occurrence), and with it the advantage over the static model.

\subsection{Practical implications}
For a SOC, the result says what to monitor and where the monitor fails. The receiving-side wave is a
property of the organization's own communication graph, computable from headers alone, so it fits an in-house
triage pipeline that does not inspect content. The blind spot is equally concrete: a single-recipient
transfer request from a compromised account is invisible to every detector studied here, so out-of-band
verification of transfers remains the control for that case, and the cascade detector should be paired with
it rather than replace it. Two further boundaries shape deployment. The learned model is the most sensitive to
noisy training labels, so retraining on analyst-confirmed labels matters more for it than for the volumetric
baseline; and at a realistic base rate every detected attack comes with hundreds of false alarms, so the
detectors are ranking instruments for analysts, not autonomous blockers. Combining header-only scoring on
the graph with content and process controls is the realistic operating mode.

\subsection{Limitations}
The limitations define the scope of these claims. First, the maliciousness is synthetic. No public corpus
links real accounts with labeled compromises and their propagation, so we generate cascades on real structure
(Enron, SNAP) or on a synthetic population. We validate the mechanism (a paced cascade separates from
legitimate traffic by its receiving-side dynamics), not the detection of real intrusions; the real-graph
replications establish the receiving-side signal on real topologies and rhythms, but the injected cascades
follow our threat model. Second, the cascade parameters ($p_\mathrm{inf}$, depth, fan-out, the stealthy
attacker's pacing) are ours; the sweeps in the Supplement show the hierarchy is stable across them, but an
attacker who sends to a single recipient (fan-out~$1$) is outside the signal by construction. Third, the scale is limited
($N{=}150$ accounts in the case study; real graphs up to about 2k nodes), making the comparison with heavy
temporal architectures a sanity check rather than a full benchmark. The case-study graph is also denser than
typical corporate mail: every account corresponds with everyone in its firm and only $7\%$ of the contact
edges cross firms, so the synthetic cascades meet a compact, well-connected neighborhood. The replications on
Enron and the SNAP graphs, whose density and rhythms are real, are what carries the mechanism beyond that
setting; the SNAP replication rests on $3$~seeds, so we read it as a regime-level confirmation rather than as
a precise effect size. Finally, the operational threshold remains
a gap. Even the best operating point (sensitivity@FPR$1\%$) is orders of magnitude from the deployment
threshold on real mail (Ho et al.\ \citep{ho2019detecting}: $<\!4$~FP/M); we report that gap as the open
requirement, not as a solved one.

\subsection{Future work directions}
The most important next step is validation on a corpus with labeled, real compromises, turning the proof of
mechanism into a measurement of detectability. The threat model should then be stressed: an attacker that
\emph{learns} the communication graph and tunes both the pacing and the propagation pattern to stay below the
receiving-side signal, together with population drift over time (TESSERACT \citep{pendlebury2019tesseract}).
On the modeling side, richer relational context---the organization's hierarchy, competencies, and public
footprint as additional layers over the same accounts---is the natural complement for the external
impersonation mode, which we study separately. Finally, the receiving-side score is a natural input to a
SOC's triage queue, provided an evaluation ties it to a real reduction in analyst load at a fixed alert
budget.
